\documentclass[../main.tex]{subfiles}
\graphicspath{{\subfix{../img/}}}
\begin{document}
	\subsection{Incipit}
	L' osservazione di fenomeni di natura elettrica risale al settimo secolo a.C, quando si scopri che l'ambra, l'ebanite e altri materiali, strofinati con un panno di lana, acquistano la proprietà di attirare corpuscoli leggeri, quali granelli di polvere e pagliuzze.\\
	William Gilbert chiamò elettrizzati i materiali che acquistano la proprietà di attirare i corpuscoli leggeri e forza elettrica la forza che si manifesta (dal termine electron che è il nome greco dell'ambra).\\
	Oggi noi attribuiamo queste forze ad interazioni tra cariche elettriche che esistono nei corpi e che passano da un corpo all'altro durante lo strofinio, per cui i corpi elettrizzati si chiamano anche elettricamente carichi.\\\\
	L'interazione elettromagnetica tiene insieme nuclei ed elettroni per formare gli atomi, tiene insieme gli atomi per formare le molecole e tiene insieme le molecole per formare gli oggetti macroscopici.\\
	Molti dei fenomeni che vediamo verificarsi attorno a noi sono in ultima analisi effetti delle forze elettromagnetiche.\\
	Sia i fenomeni elettrici che quelli magnetici dipendono dalla stessa proprietà della materia, una proprietà che chiamiamo carica elettrica.\\\\
	Supponiamo di strofinare l'estremità dì una bacchetta di vetro con un panno di seta e di appendere la bacchetta a un filo.\\
	Poi strofiniamo, sempre con un panno di seta, l'estremità di un'altra bacchetta di vetro:
	\begin{center}
		\begin{figure}[h]
			\centering
			\includegraphics[scale=0.55]{1.png}
			\caption{}
			\label{fig:im-1}
		\end{figure}
	\end{center}
	Si dice che le bacchette sono diventate "elettricamente cariche, cioè che "posseggono carica elettrica", e la forza che esercitano l'una sull'altra è chiamata "forza elettrica".
	Se si usano bacchette di plastica e un pezzo di pelliccia:
	\begin{center}
		\begin{figure}[h]
			\centering
			\includegraphics[scale=0.55]{2.png}
			\caption{}
			\label{fig:im-2}
		\end{figure}
	\end{center}
	Se poi l'estremità carica della bacchetta di vetro viene avvicinata all'estremità carica della bacchetta di plastica sospesa
	\begin{center}
		\begin{figure}[h]
			\centering
			\includegraphics[scale=0.55]{3.png}
			\caption{}
			\label{fig:im-3}
		\end{figure}
	\end{center}
	Una bacchetta di vetro che sia stata caricata strofinandola con un panno di seta è attratta dal panno. \\
	Quindi anche il panno di seta è carico.\\
	Analogamente nell’altro caso.
	\subsection{Carica Elettrica}
	Compiendo esperimenti di questo genere con bacchette fatte di molti materiali diversi e strofinate con panni di molti tipi diversi si trova:
	\begin{enumerate}
		\item La materia contiene due tipi di carica elettrica, la carica detta positiva e quella negativa.\\
		Nei corpi elettricamente neutri ci sono quantità uguali dei due tipi di carica.\\
		Quando i corpi vengono caricati per strofinio, la carica elettrica si trasferisce dall'uno all'altro.\\
		Uno dei corpi mostra un eccesso di carica positiva e l'altro un eccesso di carica negativa.
		\item Corpi che hanno cariche in eccesso dello stesso segno si respingono.
		\item Corpi che hanno cariche in eccesso di segno opposto si attraggono.
	\end{enumerate}
	\begin{center}
		\begin{figure}[h]
			\centering
			\includegraphics[scale=0.55]{4.png}
			\caption{}
			\label{fig:im-4}
		\end{figure}
	\end{center}
	\subsubsection{Struttura Atomica}
	\begin{equation}
		\begin{cases}
			m_p = 1.6725 \cdot 10^{-27} kg\\
			m_n = 1.6748 \cdot 10^{-27} kg\\
			m_e = 9.1091 \cdot 10^{-31} kg
		\end{cases}
	\end{equation}
	Dove $m_p$, $m_n$ e $m_e$ indicano rispettivamente la massa di protoni, neutroni ed elettroni.
	\subsubsection{Conduttori e Isolanti}
	Un conduttore è un materiale attraverso il quale la carica può fluire facilmente, mentre un isolante è un materiale che non consente alla carica di fluire liberamente.\\
	I metalli sono solitamente buoni conduttori mentre i nonmetalli sono di solito buoni isolanti.\\
	Generalmente quando la carica si muove attraverso un materiale sono gli elettroni che si spostano. Si dice che gli elettroni sono i portatori di carica.\\\\
	Materiale \textbf{conduttore}: reticolo di ioni positivi fissi + elettroni "liberi" in grado di muoversi attraverso il materiale.\\
	Il numero degli elettroni liberi presenti in un conduttore dipende dal materiale, ma è dell'ordine di uno per atomo.\\\\
	In un \textbf{isolante} gli elettroni sono trattenuti dai propri atomi e non possono passare da un sito a quello contiguo.\\\\
	In un fluido conduttore, quale l'acqua salata il sale nell'acqua si dissocia in ioni positivi e negativi, e gli ioni, che sono in grado di muoversi, sono i portatori di carica.
	\subsubsection{Induzione Elettrostatica}
	\begin{center}
		\begin{figure}[h]
			\centering
			\includegraphics[scale=0.55]{5.png}
			\caption{}
			\label{fig:im-5}
		\end{figure}
	\end{center}
	Quando la bacchetta carica positivamente si trova vicino alla sfera conduttrice, la parte della sfera più vicina alla bacchetta ha un eccesso di carica negativa e la parte più lontana ha un eccesso di carica positiva perché:
	\begin{itemize}
		\item Cariche di segno opposto si attraggono e cariche dello stesso segno si respingono;
		\item La carica è in grado di muoversi agevolmente attraverso la sfera conduttrice.
	\end{itemize}
	Se la sfera è connessa a terra, le cariche positive possono allontanarsi ulteriormente dalla bacchetta carica percorrendo il filo di messa a terra.\\
	Una volta che il filo è stato rimosso la carica negativa in eccesso rimane sulla sfera.\\
	Una volta allontanata anche la bacchetta la carica in eccesso rimane (l’aria secca e’ un buon isolante). \\\\
	Abbiamo illustrato un processo di carica per induzione, perche’ il processo è avvenuto SENZA CONTATTO tra bacchetta e sfera.
	\subsubsection{Quantizzazione della Carica Elettrica}
	Le particelle elementari di cui è costituita la materia sono elettricamente neutre o hanno una carica che in modulo ha un ben
	determinato valore indicato con la lettera $e$ (–$e$ è la carica dell'elettrone):
	\begin{equation}
		e = 1.60207 \cdot 10^{-19} C
	\end{equation}
	dove $\left[ C \right]$ è il simbolo dell'unità di misura nel SI, il \textbf{Coulomb}.\\\\
	Si può quindi assumere che la carica elettrica di un qualunque oggetto sia un multiplo intero, positivo o negativo, della carica
	dell'elettrone.\\
	Il valore della carica elementare è molto piccolo e il numero di cariche coinvolte è di solito molto grande, per cui spesso si tratta la carica elettrica come una grandezza continua e non discreta.
	\subsubsection{Conservazione della Carica Elettrica}
	\begin{definition}
		La somma algebrica delle cariche elettriche si mantiene costante nel tempo.
	\end{definition}
	Nei vari processi fisici (ad es. strofinio, contatto) si producono spostamenti di cariche da un corpo ad un altro, ma non si realizza mai la creazione di cariche la cui somma algebrica sia diversa da zero.\\
	Ad esempio nello strofinio si genera una carica $q$ su un corpo e –$q$ sull’altro.
	\subsection{Legge di Coulomb}
	La legge che descrive la forza che agisce tra particelle cariche è stata determinata nel 1784 da Charles Augustin Coulomb.\\
	Servendosi di una bilancia di torsione, egli determinò la dipendenza della forza elettrica dalla distanza e dalla carica. \\
	In suo onore, l'unità SI della carica è chiamata coulomb (C).
	\begin{center}
		\begin{figure}[h]
			\centering
			\includegraphics[scale=0.45]{6.png}
			\caption{}
			\label{fig:im-6}
		\end{figure}
	\end{center}
	Consideriamo l'interazione elettrica tra due particelle che possiedono carica $q_1$ e $q_2$.\\
	L'esperimento mostra che l'espressione della forza $f_{21}$ esercitata da 1 su 2 è:
	\begin{equation}
		\vec{f}_{21}=\frac{1}{4\pi\varepsilon_{0}}\frac{q_{1}q_{2}}{r^{2}}\hat{r}_{21}
	\end{equation}
	Dove $\varepsilon_{0}$ indica la costante dielettrica del vuoto: $\varepsilon_{0}=8.85\times10^{-12} \frac{C^{2}}{(Nm^{2})}$
	\begin{equation}
		\frac{1}{4\pi\varepsilon_0}=8.987\times10^9 \frac{N~m^2}{{C}^2} \cong9\times10^9 \frac{Nm^{2}}{C^2}
	\end{equation}
	\begin{itemize}
		\item La forza di Coulomb agente tra particelle cariche è inversamente proporzionale al quadrato della distanza (come la forza gravitazionale).
		\item Nella legge di Coulomb è contenuto il fatto che particelle con cariche dello stesso segno si respingono e particelle con cariche di segno opposto si attraggono.
		\item La forza di Coulomb obbedisce alla terza legge di Newton. \\
		Se $f_{ba}$ è la forza esercitata da $a$ su $b$ e la forza $f_{ab}$ esercitata da $b$ su $a$, allora $f_{ab} = -f_{ba}$ e le due forze hanno la stessa retta d’azione.
	\end{itemize}
	\subsection{Principio di sovrapposizione}
	La forza elettrica agente su di una carica elettrica $q$ in presenza di due o più cariche elettriche $(q_{1}, q_{2}, \ldots, q_{n})$ è la somma vettoriale delle forze elettriche $(f_{1}, f_{2}, \ldots, f_{n})$ esercitate su $q$ da ciascuna delle $n$ cariche elettriche considerate singolarmente:
	$$ \boldsymbol {f} = \boldsymbol {f} _ {1} + \boldsymbol {f} _ {2} + \dots + \boldsymbol {f} _ {n} $$
	Equivale ad affermare che l’interazione tra due cariche elettriche e’ indipendente dalla presenza di altre cariche
	\subsection{Confronto tra forza gravitazionale ed elettrostatica}
	Dati un protone ed un elettrone, si calcoli il rapporto tra le intensita’ delle loro forze  gravitazionali ed elettriche:
	\begin{center}
		\begin{figure}[h]
			\centering
			\includegraphics[scale=0.55]{7.png}
			\caption{}
			\label{fig:im-7}
		\end{figure}
	\end{center}
	Ricordiamo che la carica di un elettrone vale: $e = 1.6 \times 10^{-19} C$\\\\
	Sappiamo che la massa di un elettrone vale: $m_e = 0.91 \times 10^{-30} Kg$\\
	Mentre la massa di un protone vale: $m_p = 1849 \cdot m_e$
	$$ r_0 \simeq  0.5 \times 10^{-10}  \, m $$\\
	Dove $r_0$ indica il \textbf{raggio di Bohr}, ovvero la distanza dell'elettrone che orbita più vicino al suo nucleo.\\
	\textbf{N.B}: Quello di Bohr è un modello che vede gli elettroni coe oggetti orgitanti, ma in realtà non è propriamente corretto vedere gli elettroni come oggetti orbitanti in una determinata posizione, piuttosto come “nuvola”\\\\
	Forza elettrica: $\left| q_1 \right| = \left| q_2 \right| = e$ $\rightarrow$ $ F_E = \frac{1}{4\pi \varepsilon_0} \cdot \frac{e \cdot e}{r_0^2} $ \\
	Forza gravitazionale $ F_G = G \frac{m_p m_e}{r_0^2} $
	$$ \frac{F_E}{F_G} = \frac{\frac{1}{4\pi \varepsilon_0} \frac{e \cdot e}{r_0^2}}{G \frac{m_p m_e}{r_0^2}} $$
	$$= \frac{\frac{1}{4\pi \varepsilon_0} \frac{e \cdot e}{\bcancel{r_0^2}}}{G \frac{m_p m_e}{\bcancel{r_0^2}}}$$
	$$= \frac{1}{4\pi \varepsilon_0 G} \frac{e^2}{1840 \cdot m_e^2} $$
	Ricordiamo che $G = 6.7 \times 10^{-11} \frac{N \cdot m^2}{\text{kg}^2} $
	$$ = \frac{9 \times 10^9}{6.7 \times 10^{-11}} \times \frac{(1.6 \times 10^{-19})^2}{1.84 \times 10^3 \times (0.91 \times 10^{-30})^2}$$ 
	$$ = 2.25 \times 10^{39} $$
	Questo dimostra che è più forte la forza elettrica rispetto a quella gravitazionale.\\
	Tra un protone e un elettrone, la forza di attrazione elettrica è circa $10^{39}$ volte più intensa rispetto alla forza di attrazione gravitazionale.
	\subsection{Campo Elettrico}
	\subsubsection{Campi scalari e campi vettoriali}
	\begin{itemize}
		\item \textbf{Campo scalare}: ad ogni punto dello spazio e’ associato un numero con unita’ di misura (es: previsioni temperatura).
		\item \textbf{Campo vettoriale}: ad ogni punto dello spazio e’ associato un vettore (modulo, direzione e verso) (es: previsioni della velocita’ del vento). 
	\end{itemize}
	\subsubsection{Campo elettrico di una carica puntiforme}
	\begin{center}
		\begin{figure}[h]
			\centering
			\includegraphics[scale=0.55]{8.png}
			\caption{}
			\label{fig:im-8}
		\end{figure}
	\end{center}
	La presenza della carica $Q$ modifica le proprietà dello spazio circostante (assumiamo che $Q$ sia fissata in un punto).\\
	N.B: Dato che ogni punto può subire una forza, allora la carica $Q$ genera un campo di forze\\
	Infatti una carica $q$ (detta carica di prova) in un punto a distanza $r$ subisce una forza
	$$ \vec{F} = \frac{1}{4\pi\varepsilon_0} \frac{qQ}{r^3} \vec{r} $$
	In analogia con quanto fatto per il campo gravitazionale, definiamo campo elettrico $E$
	$$ \vec{E}(\vec{r}) \equiv \vec{E}(x,y,z) = \frac{\vec{F}}{q} $$
	Quindi possiamo dire che:
	\begin{equation}
		\vec{E}_0 (\vec{r}) = \frac{1}{4\pi\varepsilon_0} \frac{Q}{r^3} \vec{r}
	\end{equation}
	Dove $\vec{E}_0$ è un campo elettrico nel vuoto.\\\\
	$E$ ha le dimensioni di una forza diviso una carica e si misura in Newton/Coulomb o in Volt/Metro.\\
	\subsubsection{Campo elettrico di un gruppo di cariche puntiformi}
	Se ho un gruppo di cariche che generano il campo elettrico $E$, in un punto $P$ il campo è sempre definito come rapporto tra la forza elettrica $F$ esercitata dal gruppo sulla carica di prova $q$ e la carica $q$ stessa:
	$$ \vec{E}(\vec{r}) \equiv \vec{E}(x, y, z) = \frac{\vec{F}}{q} $$
	Siccome il principio di sovrapposizione vale per la forza di Coulomb, si avrà un principio di sovrapposizione per il campo elettrico.
	\begin{center}
		\begin{figure}[h]
			\centering
			\includegraphics[scale=0.55]{9.png}
			\caption{}
			\label{fig:im-9}
		\end{figure}
	\end{center}
	$$ \vec{E}_0(\vec{r}) = \frac{1}{4\pi \varepsilon_0} \frac{Q_1}{|\vec{r} - \vec{r}_1|^3} (\vec{r} - \vec{r}_1) + \frac{1}{4\pi \varepsilon_0} \frac{Q_2}{|\vec{r} - \vec{r}_2|^3} (\vec{r} - \vec{r}_2) $$
	$$ \boxed{\vec{E}_0(\vec{r}) = \frac{1}{4\pi \varepsilon_0} \sum_{i=1}^{n} \frac{Q_i}{|\vec{r} - \vec{r}_i|^3} (\vec{r} - \vec{r}_i)} $$
	Ad es.
	$$ E_{0 x}(x, y, z) = \frac{1}{4\pi \varepsilon_0} \sum_{i=1}^{n} \frac{Q_i(x - x_i)}{\left[(x - x_i)^2 + (y - y_i)^2 + (z - z_i)^2\right]^{\frac{3}{2}}} $$
	Quando si misura un campo elettrico, la carica di prova deve essere sufficientemente piccola da non alterare in modo significativo la distribuzione di carica che genera il campo.
	\subsubsection{Linee di forza di un campo elettrico (Visualizzazione dei campi vettoriali)}
	Le linee di forza del campo elettrico sono quelle linee orientate la cui tangente in ogni punto ha la stessa direzione e verso del campo $E$ in quel punto.\\
	In una data rappresentazione, la densità di linee di forza dipende dal modulo del campo. \\
	Nelle regioni in cui le linee sono fitte $E$ è grande, mentre dove sono rade $E$ è piccolo.\\
	Un campo uniforme è rappresentato da linee di forza equidistanti, rettilinee e parallele.
	$$ \vec{E}_0(\vec{r}) = \frac{1}{4\pi \epsilon_0} \frac{Q}{r^3} \vec{r} $$
	\begin{figure}[h]
		
		\begin{subfigure}{0.5\textwidth}
			\includegraphics[width=0.9\linewidth, height=6cm]{10.png} 
			\caption{$Q>0$}
			\label{fig:im-10}
		\end{subfigure}
		\begin{subfigure}{0.5\textwidth}
			\includegraphics[width=0.9\linewidth, height=6cm]{11.png}
			\caption{$Q<0$}
			\label{fig:im-11}
		\end{subfigure}
		
		\caption{Caption}
		\label{fig:image2}
	\end{figure}
	\subsection{Dipolo Elettrico}
	Un dipolo elettrico è costituito da due cariche opposte $(q_{1} = Q = -q_{2})$ a distanza fissa $d$ (in genere molto piccola rispetto alle altre distanze).
	\begin{center}
		\begin{figure}[h]
			\centering
			\includegraphics[scale=0.55]{12.png}
			\caption{Momento elettrico del dipolo}
			\label{fig:im-12}
		\end{figure}
	\end{center}
	Il \textbf{momento elettrico del dipolo} è definito come: $\vec{p} = Q \cdot \vec{d}$\\\\
	Il campo $E$ si calcola sfruttando il principio di sovrapposizione. \\
	Si trova che il campo elettrico prodotto da un dipolo elettrico in un punto a distanza grande rispetto alle dimensioni del dipolo stesso $(r \gg d)$ è
	$$ \vec {\mathbf {E}} (\vec {\mathbf {r}}) = \frac {1}{4 \pi \varepsilon_ {0}} \frac {p}{r ^ {3}} \big [ 3 (\hat {\mathbf {r}} \cdot \hat {\mathbf {p}}) \hat {\mathbf {r}} - \hat {\mathbf {p}} \big ] $$
	\begin{center}
		\begin{figure}[h]
			\centering
			\includegraphics[scale=0.55]{13.png}
			\caption{}
			\label{fig:im-13}
		\end{figure}
	\end{center}
	$p$ caratterizza completamente il dipolo: conoscendone il valore è possibile calcolare gli effetti del dipolo sull’ambiente circostante e gli effetti dell’ambiente circostante sul dipolo.
	\subsubsection{Calcolo campo elettrico nel campo equatoriale del dipolo}
	\begin{center}
		\begin{figure}[h]
			\centering
			\includegraphics[scale=0.55]{14.png}
			\caption{Calcolo campo elettrico nel campo equatoriale del dipolo}
			\label{fig:im-14}
		\end{figure}
	\end{center}
	$$\vec{r} \equiv (x,y,z) = (0,D,0)$$
	$$\vec{r}_1 \equiv (x_1,y_1,z_1) = (0,0,\frac{d}{2})$$
	$$\vec{r}_2 \equiv (x_2,y_2,z_2) = (0,0,-\frac{d}{2})$$
	Distanza tra punto e carica (Norma del vettore):
	$$ \left| \vec{r} - \vec{r_1} \right| = \sqrt{(x - x_1)^2 + (y - y_1)^2 + (z - z_1)^2} = \sqrt{D^2 + \left(\frac{d}{2}\right)^2} \equiv \Delta $$
	Campo elettrico all’origine in direzione $x$, ha senso che sia zero
	$$ E_{0x} = \frac{1}{4\pi \varepsilon_0} \frac{1}{\Delta^3} \left[ Q(x - x_1) - Q(x - x_2) \right] = 0 $$
	Campo elettrico all’origine in direzione $y$.\\
	I contributi delle due cariche si annullano su entrambe gli assi.
	$$ E_{0y} = \frac{1}{4\pi \varepsilon_0} \frac{1}{\Delta^3} \left[ Q(y - y_1) - Q(y - y_2) \right] = \frac{1}{4\pi \varepsilon_0} \frac{1}{\Delta^3} [QD - QD] = 0 $$
	$\uparrow$Stessa cosa, il dipolo non può esercitare forza sull’asse y quando la particella è nell’origine
	$$ E_{0z} = \frac{1}{4\pi \varepsilon_0} \frac{1}{\Delta^3} \left[ Q(z - z_1) - Q(z - z_2) \right] $$
	$$= \frac{1}{4\pi \varepsilon_0} \frac{1}{\Delta^3} \left[ Q\left(-\frac{d}{2}\right) - Q \left(\frac{d}{2}\right) \right] = \frac{1}{4\pi \varepsilon_0} \frac{-Qd}{\Delta^3} $$
	$$\vec{E}_0(0, D, 0) = \frac{-\vec{P}}{4\pi \varepsilon_0 \left[ D^2 + \left(\frac{d}{2}\right)^2 \right]^{3/2}}$$
	\subsection{Campo elettrico generato da distribuzioni continue di carica}
	La carica che si manifesta nei corpi macroscopici, come le bacchette di prima, è dovuta a uno squilibrio tra il numero di elettroni e di protoni.\\
	Le cariche presenti su corpi macroscopici sono dovute a grandi numeri di elettroni in eccesso o in difetto. \\
	Quindi è possibile trattare la carica come una distribuzione continua di elementi infinitesimi di carica $dq$.
	Densità spaziale di carica $\rho(x,y,z)$
	$$ dq = \rho(x, y, z) \, d\tau $$
	Dove $d\tau = dx \; dy \; dz$
	Integriamo una densità di carica in un volume
	$$ \vec{E}_0(\vec{r}) = \frac{1}{4\pi\varepsilon_0} \int_\tau \frac{dq(\vec{r'})}{\left| \vec{r} - \vec{r}' \right|^3} (\vec{r} - \vec{r}') $$
	$$ \boxed{\vec{E}_0(\vec{r}) = \frac{1}{4\pi\varepsilon_0} \int_\tau \frac{\rho(x', y', z')}{\left| \vec{r} - \vec{r}' \right|^3} (\vec{r} - \vec{r}') , dx' , dy' , dz'} $$
	\begin{center}
		\begin{figure}[h]
			\centering
			\includegraphics[scale=0.55]{15.png}
			\caption{Densità superficiale di carica $\sigma(x,y,z)$}
			\label{fig:im-15}
		\end{figure}
	\end{center}
	$$ dq = \sigma(x, y, z) \, ds $$
	$$ \vec{E_0}(\vec{r}) = \frac{1}{4\pi \varepsilon_0} \int_{\Sigma} \frac{\sigma(x', y', z') (\vec{r} - \vec{r}')}{\left| \vec{r} - \vec{r}' \right|^3} \, ds' $$
	\begin{center}
		\begin{figure}[h]
			\centering
			\includegraphics[scale=0.55]{16.png}
			\caption{Densità lineare di carica $\lambda(x,y,z)$}
			\label{fig:im-16}
		\end{figure}
	\end{center}
	$$ dq = \lambda(x, y, z) \, dl $$
	$$ \vec{E_0}(\vec{r}) = \frac{1}{4\pi \varepsilon_0} \int_{\Lambda} \frac{\lambda(x', y', z') (\vec{r} - \vec{r}')}{\left| \vec{r} - \vec{r}' \right|^3} \, dl’ $$
	\subsection{Campo elettrico di una distribuzione lineare di carica}
	Vogliamo calcolare il campo nel punto $P$.
	\begin{center}
		\begin{figure}[h]
			\centering
			\includegraphics[scale=0.55]{17.png}
			\caption{}
			\label{fig:im-17}
		\end{figure}
	\end{center}
	$$ \lambda = \frac{Q}{L}$$
	$$ r = \frac {R}{\cos \theta} $$
	$$ l = R \tan \theta $$
	$$ d l = R \frac {d \theta}{\cos^ {2} \theta} $$
	Prendo due punti simmetrici per poi integrare da $0$ a infinito invece che da -inf a +inf.
	$$ d\vec {E _ {0}} = d\vec { E _ {1}} + d\vec { E _ {2}} = ( d E _ {1} + d E _ {2}) \cos \theta \hat {n} \ $$
	$$ = \frac {2 \cos \theta}{4 \pi \varepsilon_ {0}} \frac {\lambda d l}{r ^ {2}} \hat {n} $$
	$$\qquad$$
	$$ \vec {E _ {0}} = \int d \vec {E _ {0}} = \frac {\lambda \hat {n}}{2 \pi \varepsilon_ {0}} \int_ {0} ^ {\infty} \frac {\cos \theta}{r ^ {2}} d l $$
	Espando $dl$ e $r$, facendo riferimento all'iizio del paragrafo dove sono definite
	$$ = \frac {\lambda \hat {n}}{2 \pi \varepsilon_ {0}} \int_ {0} ^ {\pi / 2} \frac {\cos^ {\cancel{3}} \theta}{R ^ {\cancel{2}}} \times \cancel{R} \frac { d \theta}{\cancel{\cos^ {2} \theta}} $$
	$$ = \frac {\lambda}{2 \pi \varepsilon_ {0} R} \hat {n} \int_ {0} ^ {\pi / 2} \cos \theta d \theta $$
	$$ \vec {E _ {0}} = \frac {1}{2 \pi \varepsilon_ {0}} \frac {\lambda}{R} \hat {n} $$
	\subsection{Spira circolare uniformemente carica (campo lungo l’asse)}
	\begin{center}
		\begin{figure}[h]
			\centering
			\includegraphics[scale=0.55]{18.png}
			\caption{}
			\label{fig:im-18}
		\end{figure}
	\end{center}
	$$ \lambda = \frac{Q}{2\pi R} $$
	Calcolo il micro-campo diretto sul punto per una lunghezza infinitesima del filo
	$$ \vec{dE_0} = \frac{1}{4\pi \varepsilon_0} \frac{\lambda dl}{r^3} \vec{r} $$
	Calcolo la componente tangente all’asse del campo infinitesimo
	$$ dE_{0x} = \frac{1}{4\pi \varepsilon_0} \frac{\lambda dl}{r^2} \cos \alpha $$
	Integrazione di tutti i campi infinitesimi
	$$ E_{0x} = \int dE_{0x} $$
	Espando $dE_{0x}$:
	$$ = \frac{1}{4\pi \varepsilon_0} \int_{\text{Spira}} \frac{\lambda \cos \alpha}{r^2} dl = \frac{\lambda \cos \alpha}{4\pi \varepsilon_0 r^2} \int_{\text{Spira}} dl = \frac{2\pi R}{4\pi \varepsilon_0} \frac{\lambda \cos \alpha}{r^2} $$
	$$ = \frac{\lambda R}{2\varepsilon_0} \frac{\cos \alpha}{r^2} = \frac{\lambda R}{2\varepsilon_0} \frac{x}{r^3} = \frac{\lambda R}{2\varepsilon_0} \frac{x}{(x^2 + R^2)^{\frac{3}{2}}} $$
	$$ E_{0x}(0) = 0 $$
	$$ E_{0x}(x) \xrightarrow[x \to \infty]{} 0 $$
	\subsection{Flusso di un campo vettoriale}
	Il concetto è stato originariamente introdotto nella teoria dei fluidi, dove il flusso è legato alla quantità di fluido che passa attraverso una data superficie geometrica, quale per esempio la sezione di un condotto.\\
	Il flusso $\Phi$ di un campo vettoriale è una grandezza scalare che dipende dal campo e dalla superficie rispetto alla quale viene calcolato.\\
	\begin{center}
		\begin{figure}[h]
			\centering
			\includegraphics[scale=0.55]{19.png}
			\caption{Campo vettoriale A uniforme sulla superficie}
			\label{fig:im-19}
		\end{figure}
	\end{center}
	Il vettore superficie lo chiamiamo $\delta S$ perchè dopo verrà usato come $dS$ in un integrale\\\\
	Per superfici chiuse, convenzione vettore normale uscente.\\
	Si introduce un vettore superficie $\Delta S$ con modulo uguale all’area della superficie e direzione perpendicolare alla superficie stessa (due versi possibili). 
	$$ \Delta \vec{S} = \Delta S \hat{n} $$
	Definizione di flusso: $ \boxed{\Phi = \vec{A} \cdot \Delta \vec{S}} $
	\begin{center}
		\begin{figure}[h]
			\centering
			\includegraphics[scale=0.55]{20.png}
			\caption{}
			\label{fig:im-20}
		\end{figure}
	\end{center}
	$$ \boxed{\Phi = \vec {A} \cdot \Delta \vec {S} = A \cdot \Delta S \cdot \cos \alpha} $$
	$\Phi$ è definito come un prodotto scalare, quindi se i vettori sono ortogonali il flusso è 0.\\\\
	Per farsi un’idea intuitiva del flusso si può ricorrere alle linee di forza: il numero delle linee che attraversano una superficie è proporzionale al flusso relativo a tale superficie.
	\begin{figure}[h]
		
		\begin{subfigure}{0.5\textwidth}
			\includegraphics[scale=0.55]{21.png} 
			\caption{}
			\label{fig:im-21}
		\end{subfigure}
		\begin{subfigure}{0.5\textwidth}
			\includegraphics[scale=0.55]{22.png}
			\caption{}
			\label{fig:im-22}
		\end{subfigure}
		
		\caption{}
		\label{fig:image3}
	\end{figure}
	L’ambiguità nel verso della normale alla superficie può essere eliminata quando la superficie è chiusa: si considera la normale uscente.\\\\
	Se la superficie $S$ è curva o il campo elettrico varia da punto a punto, il flusso si calcola dividendo $S$ in piccoli elementi di superficie, ciascuno abbastanza piccolo da poter essere considerato piano e tale che su di esso la variazione del campo elettrico possa essere trascurata. \\
	Il flusso attraverso l’intera superficie è allora la somma dei contributi dovuti a ciascuno dei piccoli elementi di superficie.\\\\
	Facendo tendere a zero le dimensioni di ciascun elemento ed a infinito il loro numero, la somma diventa un integrale:
	\begin{center}
		\begin{figure}[h]
			\centering
			\includegraphics[scale=0.55]{23.png}
			\caption{}
			\label{fig:im-23}
		\end{figure}
	\end{center}
	$$ \vec {d s} = \hat {n} \; d s $$
	$$ d \phi (\vec {A}) = \vec {A} \cdot \vec {d s} = \vec {A} \cdot \hat {n} \; d s = A \; d s \cos \theta $$
	$$ \Phi_ {S} (\vec {A}) \equiv \int_ {S} d \phi (\vec {A}) = \int_ {S} \vec {A} \cdot \vec {d s} $$
	Dove notiamo che $\Phi_ {S} (\vec {A})$ equivale ad un integrale in $S$ si un prodotto scalare.
	\subsection{Legge di Gauss}
	\begin{definition}
			Il flusso del campo elettrico attraverso una qualunque superficie chiusa (come una sfera: una superficie in 3 dimensioni chiusa) è pari alla somma algebrica delle cariche contenute all’interno del volume racchiuso dalla superficie divisa per la costante dielettrica del vuoto.
	\end{definition}
	$$ \Phi_ {S} (\vec {E _ {0}}) \equiv \int_ {S} \vec {E _ {0}} \cdot \hat {\vec {n}} \; d s = \frac {1}{\varepsilon _ {o}} \sum Q _ {int} = \frac {Q ^{int}_ {tot}}{\varepsilon _ {0}} $$
	Dove $\hat{n}$ è la normale.\\\\
	La superficie chiusa attraverso la quale si calcola il flusso del campo elettrico generalmente è una superficie geometrica immaginaria che non corrisponde alla superficie di un oggetto.\\
	Il campo vettoriale esterno in una superficie chiusa non provoca alcun flusso, quindi rimane solo quello provocato dalle cariche interne.
	\begin{center}
		\begin{figure}[h]
			\centering
			\includegraphics[scale=0.55]{24.png}
			\caption{}
			\label{fig:im-24}
		\end{figure}
	\end{center}
	Il campo E nell’integrale del flusso è il campo dovuto a tutte le particelle cariche, sia interne che esterne alla superficie considerata, ma il flusso attraverso l’intera superficie è dovuto soltanto alle cariche che si trovano all’interno.
	\subsection{Potenziale Elettrico}
	\subsection{Energia Potenziale Elettrica}
	\subsection{Energia Elettrostatica}
	\subsection{Condensatori}
\end{document}